% GENERATED FILE. Edit canonical lessons, then run npm run generate:books.
% canonical-source-hash: fnv1a64:e5fc8a289d19656d
% canonical-lessons: ES-C768-hervir, ES-C768-freir, ES-C768-pelar, ES-C768-batir, ES-C768-ingrediente

\chapter{Lo que cociné}
\label{ch:es-lo-que-cocine}
\hlchaptermodality{768}

\begin{chapteropening}
\textbf{By the end of this chapter:} \emph{I can tell someone how I cooked a dish: what I boiled, fried, peeled and beat, and which ingredients I used.}

Complete the last lesson of chapter 768: I can tell someone how I cooked a dish: what I boiled, fried, peeled and beat, and which ingredients I used.
\end{chapteropening}

\section[hervir]{hervir — to boil}

\label{lesson:ES-C768-hervir}



\begin{quote}
Before the new one: say the Spanish for to seem, then the Spanish for to keep.
\end{quote}

\subsection*{You'll want to know: hervir}
\textbf{hervir} — \textquotedblleft{}to boil\textquotedblright{}.

\begin{quote}
Herví el agua para la pasta.
\end{quote}

\begin{quote}
\emph{I boiled the water for the pasta.}
\end{quote}

Stem change: \emph{hiervo}, \emph{hirvió}. \emph{El agua está hirviendo} — the water is boiling.

\begin{cousinweb}
\emph{hervir} is inherited from Latin \emph{fervēre} 'to boil, to seethe' (EtymDB), through old Spanish \emph{herver}. Latin \emph{f-} became Spanish \emph{h-}. English \emph{fervent} comes from the same verb.
\end{cousinweb}

\begin{grammarlens}[title={What you've built}]
One more word for telling, explaining and describing at B1.
\end{grammarlens}

\subsection*{Your turn}
\begin{itemize}
\raggedright
  \item \emph{Say it:} \emph{hervir}
  \item \emph{Say it:} \emph{hervir}, once more
  \item \emph{Say it:} a sentence of your own with \emph{hervir}
  \item \emph{Recall:} say the Spanish for to seem, then the Spanish for to keep, then say \emph{hervir} again
\end{itemize}

\begin{tcolorbox}[breakable,colback=teal!4,colframe=teal!35!black,title={Before you move on}]
What does hervir mean? (\textquotedblleft{}To boil\textquotedblright{}.) Say it once more.
\end{tcolorbox}
\section[freír]{freír — to fry}

\label{lesson:ES-C768-freir}



\begin{quote}
Before the new one: say the Spanish for to keep, then the Spanish for to boil.
\end{quote}

\subsection*{You'll want to know: freír}
\textbf{freír} — \textquotedblleft{}to fry\textquotedblright{}.

\begin{quote}
Frió los huevos con un poco de aceite.
\end{quote}

\begin{quote}
\emph{He fried the eggs in a little oil.}
\end{quote}

Stem change: \emph{frío}, \emph{frió}. Participle: \emph{frito} — \emph{patatas fritas}.

\begin{cousinweb}
\emph{freír} is inherited from Latin \emph{frīgere} 'to roast, to fry' (EtymDB). The \emph{g} between vowels was lost. It is not the \emph{frigidus} 'cold' behind \emph{frío}.
\end{cousinweb}

\begin{grammarlens}[title={What you've built}]
One more word for telling, explaining and describing at B1.
\end{grammarlens}

\subsection*{Your turn}
\begin{itemize}
\raggedright
  \item \emph{Say it:} \emph{freír}
  \item \emph{Say it:} \emph{freír}, once more
  \item \emph{Say it:} a sentence of your own with \emph{freír}
  \item \emph{Recall:} say the Spanish for to keep, then the Spanish for to boil, then say \emph{freír} again
\end{itemize}

\begin{tcolorbox}[breakable,colback=teal!4,colframe=teal!35!black,title={Before you move on}]
What does freír mean? (\textquotedblleft{}To fry\textquotedblright{}.) Say it once more.
\end{tcolorbox}
\section[pelar]{pelar — to peel}

\label{lesson:ES-C768-pelar}



\begin{quote}
Before the new one: say the Spanish for to boil, then the Spanish for to fry.
\end{quote}

\subsection*{You'll want to know: pelar}
\textbf{pelar} — \textquotedblleft{}to peel\textquotedblright{}.

\begin{quote}
Pelé las patatas y las corté.
\end{quote}

\begin{quote}
\emph{I peeled the potatoes and cut them.}
\end{quote}

\emph{Pelarse} is to peel after sunburn. \emph{Hace un frío que pela} — it is bitterly cold.

\begin{cousinweb}
\emph{pelar} is inherited from Latin \emph{pilāre} 'to strip of hair' (EtymDB), built on \emph{pilus} 'a hair', the word behind \emph{pelo}. Peeling a fruit is stripping its skin.
\end{cousinweb}

\begin{grammarlens}[title={What you've built}]
One more word for telling, explaining and describing at B1.
\end{grammarlens}

\subsection*{Your turn}
\begin{itemize}
\raggedright
  \item \emph{Say it:} \emph{pelar}
  \item \emph{Say it:} \emph{pelar}, once more
  \item \emph{Say it:} a sentence of your own with \emph{pelar}
  \item \emph{Recall:} say the Spanish for to boil, then the Spanish for to fry, then say \emph{pelar} again
\end{itemize}

\begin{tcolorbox}[breakable,colback=teal!4,colframe=teal!35!black,title={Before you move on}]
What does pelar mean? (\textquotedblleft{}To peel\textquotedblright{}.) Say it once more.
\end{tcolorbox}
\section[batir]{batir — to beat, to whisk}

\label{lesson:ES-C768-batir}



\begin{quote}
Before the new one: say the Spanish for to fry, then the Spanish for to peel.
\end{quote}

\subsection*{You'll want to know: batir}
\textbf{batir} — \textquotedblleft{}to beat, to whisk\textquotedblright{}.

\begin{quote}
Batí tres huevos para la tortilla.
\end{quote}

\begin{quote}
\emph{I beat three eggs for the omelette.}
\end{quote}

\emph{Batir un récord} — to break a record.

\begin{cousinweb}
\emph{batir} comes from Late Latin \emph{battere} 'to beat' (EtymDB), a form of Latin \emph{battuere}. \emph{Batido} 'milkshake' is its participle.
\end{cousinweb}

\begin{grammarlens}[title={What you've built}]
One more word for telling, explaining and describing at B1.
\end{grammarlens}

\subsection*{Your turn}
\begin{itemize}
\raggedright
  \item \emph{Say it:} \emph{batir}
  \item \emph{Say it:} \emph{batir}, once more
  \item \emph{Say it:} a sentence of your own with \emph{batir}
  \item \emph{Recall:} say the Spanish for to fry, then the Spanish for to peel, then say \emph{batir} again
\end{itemize}

\begin{tcolorbox}[breakable,colback=teal!4,colframe=teal!35!black,title={Before you move on}]
What does batir mean? (\textquotedblleft{}To beat, to whisk\textquotedblright{}.) Say it once more.
\end{tcolorbox}
\section[el ingrediente]{el ingrediente — the ingredient}

\label{lesson:ES-C768-ingrediente}



\begin{quote}
Before the new one: say the Spanish for to peel, then the Spanish for to beat.
\end{quote}

\subsection*{You'll want to know: el ingrediente}
\textbf{el ingrediente} — \textquotedblleft{}the ingredient\textquotedblright{}.

\begin{quote}
Usé solo cuatro ingredientes para la salsa.
\end{quote}

\begin{quote}
\emph{I used only four ingredients for the sauce.}
\end{quote}

\begin{cousinweb}
\emph{ingrediente} is a learned borrowing from Latin \emph{ingrediēns} 'going in' (EtymDB), from \emph{ingredī} 'to enter', built on \emph{in} and \emph{gradī} 'to step'. An ingredient is what goes into the dish.
\end{cousinweb}

\begin{grammarlens}[title={What you've built}]
One more word for telling, explaining and describing at B1.
\end{grammarlens}

\subsection*{Your turn}
\begin{itemize}
\raggedright
  \item \emph{Say it:} \emph{el ingrediente}
  \item \emph{Say it:} \emph{el ingrediente}, once more
  \item \emph{Say it:} a sentence of your own with \emph{el ingrediente}
  \item \emph{Recall:} say the Spanish for to peel, then the Spanish for to beat, then say \emph{el ingrediente} again
\end{itemize}

\begin{tcolorbox}[breakable,colback=teal!4,colframe=teal!35!black,title={Before you move on}]
What does el ingrediente mean? (\textquotedblleft{}The ingredient\textquotedblright{}.) Say it once more.
\end{tcolorbox}
